% ============================================================================
% Fichier  : partie2/P02_C03_iso31000_iso27005.tex
% Titre    : Gérer les Risques : ISO 31000 et ISO/IEC 27005
% Auteur   : MINKA MI NGUIDJOÏ Thierry Emmanuel
% Version  : 1.0 -- novembre 2026
% Corrigés : corriges/Corr_C03.tex
% ============================================================================
% Faits du cas (bible des cas) : BANQUE SANAGA, R-AG-01 (inhérent 16,
% résiduel 12), convention de réduction par étape, appétence provisoire de
% novembre 2026, registre 2025 DC-01 à DC-06, DIGICORE sans clause d'audit,
% pare-feu d'accès au siège (journaux 30 jours). Ex. 2.2 : en mars 2027,
% second facteur testé, cloisonnement non testé, sauvegardes restaurées.
% Faits nouveaux (reportés dans la bible) : critères d'acceptation de
% décembre 2026 ; identification par les deux approches (GAB sous système
% obsolète, droits d'administration locale des chefs d'agence) ; plan de
% traitement A1 à A10 (193 M FCFA pour R-AG-01) ; état en septembre 2027 :
% 5 OUI / 2 NON / 3 NEK ; décision d'acceptation de R-DC-02.
% Sources : ISO 31000:2018 (principes, cadre, processus, options de
% traitement : paraphrasés) ; ISO/IEC 27005:2022 (deux approches, critères :
% paraphrasés, sans numéro de paragraphe) ; ISO/IEC 27001:2022, art. 6.1.3
% (paraphrasé) ; loi n° 2010/012, art. 26 al. 1 et 29 al. 1 (texte français
% officiel, Ann. D d'ACSI).
% ============================================================================

\chapter{Gérer les Risques : ISO 31000 et ISO/IEC 27005}
\label{chap:iso27005}
\chaptermark{ISO 31000 et ISO/IEC 27005}

\epigraph{%
    « Un registre est une photographie.\\
    La gestion des risques est un processus. »%
}{Principe du cours}

\niveauChapitre{Fondamental}{%
    Maîtrise du Chapitre~\ref{chap:risque}%
}

\begin{boxobjectifs}
\begin{itemize}
    \item Situer les principes, le cadre et le processus de l'ISO~31000.
    \item Établir le contexte d'une appréciation~: périmètre, parties
          intéressées, critères de risque et critères d'acceptation.
    \item Appliquer les deux approches d'identification de
          l'ISO/IEC~27005:2022 et montrer qu'elles se complètent.
    \item Choisir une option de traitement et rédiger un plan de
          traitement.
    \item Faire accepter formellement le risque résiduel, et refuser une
          acceptation fondée sur une cote \NEK{}.
    \item Suivre l'avancement d'un plan de traitement par $\mu$, $U$ et
          $H$.
\end{itemize}
\end{boxobjectifs}

\bigskip

% ============================================================================
\section*{Avant de commencer : « Donnez-nous un plan »}
% ============================================================================

Douala, décembre 2026. Le comité des risques du conseil de la BANQUE
SANAGA\footnote{Comme dans tout le cours, la BANQUE SANAGA, ses
dirigeants, ses prestataires, les dates et les montants sont fictifs~; les
normes et les textes cités sont exacts.} a reçu en novembre la cotation de
R-AG-01~: 12, intolérable. Il ne demande plus une cote~; il demande un
plan. Que fait-on, qui le fait, pour quand, à quel coût, et que restera-t-il
ensuite~?

Mme~ABENA mesure l'écart. Le registre de 2025 était une liste de six
scénarios, revue une fois par an. Il n'avait ni critères d'acceptation, ni
plan de traitement, ni suivi. La banque avait un registre~; elle n'avait
pas de processus. La RSSI propose de reprendre la démarche depuis le
début, selon l'ISO~31000 et l'ISO/IEC~27005, sur toute la banque cette
fois.

\medskip

Ce chapitre ouvre la Partie~II, consacrée aux méthodes. Il présente le
processus général de gestion des risques, puis la manière dont l'ISO/IEC
27005 l'applique à la sécurité de l'information, jusqu'à l'acceptation
formelle de ce qui reste.

\begin{boxalert}
\textbf{Violation typique.} Tenir un registre des risques sans plan de
traitement ni décision d'acceptation~: les risques sont décrits, cotés, et
laissés en l'état.

\smallskip
\textbf{Règle~:} chaque risque d'un registre aboutit à une décision
écrite~: un traitement, avec un responsable et une échéance, ou une
acceptation, signée par l'instance compétente. Un risque coté sans
décision n'est pas géré.
\end{boxalert}

\separateur

% ============================================================================
\section{L'ISO 31000 : Principes, Cadre, Processus}
\label{sec:c03-iso31000}
% ============================================================================

\subsection{Une norme générale}

L'ISO~31000:2018 donne des lignes directrices pour le management du
risque, quel qu'il soit~: financier, industriel, juridique,
cyber~\cite{iso31000_2018}. Elle n'est pas certifiable et ne prescrit pas
de méthode. Elle s'articule en trois ensembles~: des \emph{principes}, qui
disent ce qu'est un management du risque efficace~; un \emph{cadre
organisationnel}, qui l'intègre dans la gouvernance~; un
\emph{processus}, qui en décrit les étapes.

\begin{tableaucours}{L{2.6cm}L{9.6cm}}{%
    \tableauHeader{Ensemble} &
    \tableauHeader{Contenu (d'après l'ISO~31000:2018)}
}
Principes & Intégré~; structuré et global~; adapté~; inclusif~;
dynamique~; fondé sur la meilleure information disponible~; attentif aux
facteurs humains et culturels~; en amélioration continue \\
Cadre & Leadership et engagement de la direction~; intégration dans les
activités~; conception~; mise en œuvre~; évaluation~; amélioration \\
Processus & Communication et consultation~; domaine d'application,
contexte et critères~; appréciation (identification, analyse,
évaluation)~; traitement~; suivi et revue~; enregistrement et rapports \\
\end{tableaucours}
\vspace{-0.8em}
\begin{center}{\small\itshape Tableau~3.1. Les trois ensembles de
l'ISO~31000}\end{center}

Le cadre fait le lien avec le Chapitre~\ref{chap:grc}~: sans leadership de
la direction ni intégration dans la gouvernance, le processus tourne à vide.
Le principe de « meilleure information disponible » fait le lien avec le
Chapitre~\ref{chap:risque}~: la norme demande que l'information soit
disponible, et elle ajoute qu'il faut tenir compte de ses limites et de ses
incertitudes. C'est exactement ce que fait la cote de la base de preuve.

\subsection{Le processus}

\begin{figure}[H]
\centering
\begin{tikzpicture}[
    etape/.style={draw=CouleurPrimaire, thick, rounded corners,
                  fill=CouleurPrimaire!8, text width=5.2cm,
                  minimum height=0.8cm, align=center, font=\small},
    transverse/.style={draw=CouleurAccent, thick, rounded corners,
                       fill=CouleurAccent!12, text width=2.4cm,
                       minimum height=5.6cm, align=center, font=\small},
    fleche/.style={->, thick, CouleurSecondaire}]
    \node[etape] (C) at (0,2.4) {Domaine d'application, contexte, critères};
    \node[etape, minimum height=1.6cm] (A) at (0,0.6)
        {\textbf{Appréciation}\\identification\\analyse\\évaluation};
    \node[etape] (T) at (0,-1.2) {Traitement};
    \node[etape] (E) at (0,-2.4) {Enregistrement et rapports};
    \draw[fleche] (C) -- (A);
    \draw[fleche] (A) -- (T);
    \draw[fleche] (T) -- (E);
    \node[transverse] (CC) at (-4.4,0) {Communication\\et\\consultation};
    \node[transverse] (SR) at (4.4,0) {Suivi\\et\\revue};
\end{tikzpicture}
\caption{Le processus de management du risque (d'après l'ISO~31000:2018)}
\label{fig:c03-processus}
\end{figure}

Deux activités encadrent les autres et durent tout le processus~: la
\emph{communication et consultation}, qui associe les parties intéressées
à chaque étape, et le \emph{suivi et revue}, qui vérifie que les décisions
sont appliquées et que les cotes restent valables. À la BANQUE SANAGA, ce
sont précisément les deux activités qui manquaient~: le registre était
rédigé par la DSI seule, et revu une fois par an.

\subsection{L'ISO/IEC 27005, déclinaison pour la sécurité de l'information}

L'ISO/IEC~27005:2022 applique ce processus à la sécurité de l'information
et le relie aux exigences de l'ISO/IEC~27001~\cite{iso27005_2022}. Elle
n'impose pas non plus de méthode~: EBIOS Risk Manager
(Chapitre~\ref{chap:ebios}) ou une méthode maison peuvent la mettre en
œuvre. Sa version de 2022 a refondu le texte de 2018~; elle insiste sur
l'identification des risques à partir des événements, et non plus
seulement des actifs.

% ============================================================================
\section{Établir le Contexte}
\label{sec:c03-contexte}
% ============================================================================

\subsection{Périmètre et parties intéressées}

Le premier livrable fixe le périmètre de l'appréciation et les parties
intéressées. La RSSI retient cette fois toute la banque~: le siège, le
centre de données, le site de secours de Bonabéri et les 38 agences. Elle
recense les parties intéressées et ce qu'elles attendent.

\begin{tableaucours}{L{3.4cm}L{8.8cm}}{%
    \tableauHeader{Partie intéressée} &
    \tableauHeader{Attente en matière de sécurité}
}
Clients & Guichets ouverts, comptes exacts, données protégées \\
Actionnaires et conseil & Pertes contenues dans l'appétence, réputation
préservée \\
COBAC & Dispositif de maîtrise des risques couvrant les systèmes
d'information~; comptes rendus après incident \\
DIGICORE et autres prestataires & Engagements contractuels clairs, accès
encadrés \\
Agents des agences & Outils disponibles, consignes simples \\
\end{tableaucours}
\vspace{-0.8em}
\begin{center}{\small\itshape Tableau~3.2. Les parties intéressées de
l'appréciation}\end{center}

\subsection{Critères de risque, critères d'acceptation}

L'ISO/IEC~27005 demande d'établir, avant toute appréciation, deux familles
de critères~: ceux qui servent à apprécier (les échelles de vraisemblance
et d'impact) et ceux qui servent à accepter~\cite{iso27005_2022}. La
banque dispose des premiers depuis novembre~2026 (Chapitre~\ref{chap:risque},
Tableau~2.3). Il lui manque les seconds~: la règle d'appétence dit quels
risques traiter, elle ne dit pas qui peut accepter quoi.

\begin{boxdef}
\termecle{Critères d'acceptation du risque}\index{critères
d'acceptation}~: règles, fixées à l'avance par l'organe de direction, qui
disent quel niveau de risque résiduel peut être accepté, par qui, à quelles
conditions et pour combien de temps.
\end{boxdef}

En décembre 2026, le comité des risques complète l'appétence par les
critères suivants.

\begin{tableaucours}{cL{3.6cm}L{7.4cm}}{%
    \tableauHeader{Résiduel} &
    \tableauHeader{Qui peut accepter} &
    \tableauHeader{Conditions}
}
1 à 4 & Le propriétaire du risque & Décision écrite~; revue annuelle \\
5 à 8 & Le comité de sécurité, sur proposition du propriétaire & Aucune
mesure supplémentaire proportionnée~; revue semestrielle~; sinon,
traitement dans l'année \\
9 et plus & Personne, sauf le comité des risques du conseil, à titre
temporaire & Plan de traitement daté~; durée de six mois au plus,
renouvelable une fois \\
Toute cote & -- & Pas d'acceptation si une cote de vraisemblance ou
d'impact a une base \NEK{} ou \NON{}, sauf acceptation « en l'état des
connaissances » avec une date de revue \\
\end{tableaucours}
\vspace{-0.8em}
\begin{center}{\small\itshape Tableau~3.3. Critères d'acceptation de la
BANQUE SANAGA (décembre 2026)}\end{center}

\begin{boxinfo}
\textbf{Pourquoi la dernière ligne~?} Accepter un risque, c'est décider en
connaissance de cause de ne pas le traiter davantage. Si la cote repose sur
une ignorance, la décision ne porte pas sur le risque, mais sur une
supposition. La règle n'interdit pas d'accepter~; elle oblige à dire ce que
l'on accepte d'ignorer, et jusqu'à quand.
\end{boxinfo}

% ============================================================================
\section{Apprécier : Deux Approches qui se Complètent}
\label{sec:c03-apprecier}
% ============================================================================

\subsection{Identifier}

L'ISO/IEC~27005:2022 propose deux approches d'identification~: par les
\emph{événements} et par les \emph{actifs}~\cite{iso27005_2022}.

\begin{tableaucours}{L{2.6cm}L{4.8cm}L{4.8cm}}{%
    \tableauHeader{} &
    \tableauHeader{Par les événements} &
    \tableauHeader{Par les actifs}
}
Point de départ & Les sources de risque, leurs objectifs, et les
événements redoutés sur les valeurs métier & L'inventaire des biens
supports, puis leurs menaces et vulnérabilités \\
Niveau & Stratégique~: on raisonne en scénarios & Détaillé~: on raisonne
composant par composant \\
Force & Rapide, centrée sur ce qui compte pour la direction & Exhaustive,
fait apparaître des faiblesses précises \\
Faiblesse & Peut oublier un composant que personne n'a imaginé & Lourde~;
peut noyer l'essentiel dans des centaines de lignes \\
\end{tableaucours}
\vspace{-0.8em}
\begin{center}{\small\itshape Tableau~3.4. Les deux approches
d'identification}\end{center}

\begin{boxscenario}{Opération SANAGA. Ce que chaque approche a trouvé}
\textbf{Par les événements.} Un atelier réunit la direction du réseau
d'agences, la DSI, la direction des risques et la direction financière. Il
retient quatre scénarios stratégiques~: le rançongiciel en agence
(R-AG-01)~; la fraude au virement par compromission d'un poste du
siège~; la défaillance de DIGICORE~; la fuite du fichier des clients.

\medskip
\textbf{Par les actifs.} L'inventaire des biens supports des 38 agences,
mené en parallèle, fait apparaître deux faiblesses qu'aucun scénario
n'avait imaginées~:
\begin{itemize}
    \item les automates de retrait (GAB) de 14 agences fonctionnent sous un
          système d'exploitation qui ne reçoit plus de correctifs~;
    \item dans 31 agences, le chef d'agence dispose des droits
          d'administration locale sur tous les postes, pour « dépanner ».
\end{itemize}

\textbf{Analyse} (confiance élevée). Les deux approches ne se recouvrent
pas. Les scénarios ont donné les priorités de la direction~; l'inventaire
a donné des faiblesses concrètes, dont la seconde alimente directement
R-AG-01 (elle facilite la propagation). On ne choisit pas entre les deux~:
on part des événements pour savoir ce qui compte, et l'on vérifie par les
actifs que rien d'important n'échappe.
\end{boxscenario}

\subsection{Analyser et évaluer}

L'\emph{analyse} cote chaque risque~: vraisemblance, impact, criticité,
avec leur base de preuve (Chapitre~\ref{chap:risque}). L'\emph{évaluation}
compare le résultat aux critères et en tire une priorité~: traiter sans
délai, traiter dans l'année, ou proposer à l'acceptation. Les deux
faiblesses de l'inventaire ne sont pas des risques à part entière~; ce sont
des vulnérabilités qui modifient la cote d'un scénario existant, ou qui en
appellent un nouveau (« compromission d'un GAB »).

% ============================================================================
\section{Traiter}
\label{sec:c03-traiter}
% ============================================================================

\subsection{Les options}

L'ISO~31000 énumère les options de traitement~: éviter le risque en
renonçant à l'activité~; prendre ou augmenter un risque pour saisir une
opportunité~; supprimer la source de risque~; modifier la vraisemblance~;
modifier les conséquences~; partager le risque~; maintenir le risque par
une décision argumentée~\cite{iso31000_2018}. En sécurité de
l'information, on les regroupe souvent en quatre familles.

\begin{tableaucours}{L{2.4cm}L{5.0cm}L{4.8cm}}{%
    \tableauHeader{Famille} &
    \tableauHeader{Principe} &
    \tableauHeader{Limite}
}
Réduire & Modifier la vraisemblance ou les conséquences par des mesures &
Ne réduit le résiduel que si les mesures sont testées \\
Partager & Transférer une partie des conséquences financières (assurance,
contrat) & Ne transfère ni la responsabilité, ni l'atteinte à la
réputation \\
Éviter & Renoncer à l'activité ou au composant qui crée le risque &
Supprime aussi la valeur que l'activité apportait \\
Maintenir & Accepter le résiduel, selon les critères & Exige une décision
écrite et une base de preuve \\
\end{tableaucours}
\vspace{-0.8em}
\begin{center}{\small\itshape Tableau~3.5. Les quatre familles de
traitement}\end{center}

\subsection{Les options pour R-AG-01}

\begin{boxscenario}{Opération SANAGA. Choisir un traitement pour R-AG-01}
\textbf{Maintenir~:} impossible. La criticité résiduelle est de 12~; seul
le comité des risques du conseil pourrait l'accepter, pour six mois, et
avec un plan.

\medskip
\textbf{Éviter~:} supprimer le partage de fichiers commun aux agences.
Option retenue en partie~: le partage commun est remplacé par un espace
par agence. Supprimer la messagerie externe en agence, en revanche,
priverait les agences d'un outil indispensable.

\medskip
\textbf{Partager~:} une assurance contre les risques cyber, si le marché
le permet. Ordre de grandeur obtenu d'un courtier~: prime de 35~millions
FCFA par an, franchise de 100~millions, plafond de 1~milliard. Les
assureurs exigent en général, avant de couvrir, un second facteur et des
sauvegardes hors ligne testées~: l'assurance vient \emph{après} la
réduction, pas à sa place.

\medskip
\textbf{Réduire~:} couvrir les trois étapes du scénario par des mesures
testées (plan ci-dessous).

\medskip
\textbf{Analyse} (confiance moyenne). La combinaison retenue est
«~réduire, éviter en partie, puis partager~». Le résiduel visé, selon la
convention de réduction, est de vraisemblance~2 (entrée et propagation
couvertes) et d'impact~3 (récupération couverte)~: criticité~6. La
convention ne permet pas de descendre plus bas~; un résiduel de~6 relève
du comité de sécurité.
\end{boxscenario}

\subsection{Le plan de traitement}

L'ISO/IEC~27001 demande, en substance, que l'organisation choisisse ses
options de traitement, détermine les mesures nécessaires, les compare à
celles de son annexe~A, établisse une déclaration d'applicabilité, formule
un plan de traitement, et obtienne des propriétaires des risques
l'approbation de ce plan et l'acceptation des risques
résiduels~\cite{iso27001_2022}. La déclaration d'applicabilité est étudiée
au Chapitre~\ref{chap:iso27000}~; voici le plan.

\begin{tableaucours}{cL{4.4cm}cL{2.2cm}cr}{%
    \tableauHeader{Réf.} &
    \tableauHeader{Action} &
    \tableauHeader{Étape} &
    \tableauHeader{Responsable} &
    \tableauHeader{Échéance} &
    \tableauHeader{M FCFA}
}
A1 & Second facteur sur tous les comptes & Entrée & DSI & mars 2027 & 25 \\
A2 & Cloisonnement du partage, un espace par agence & Propag. & DSI &
fév. 2027 & 40 \\
A3 & Test d'intrusion interne vérifiant le cloisonnement & Propag. & RSSI
& juin 2027 & 15 \\
A4 & Sauvegardes hors ligne des serveurs d'agence, test trimestriel &
Récup. & DSI & mars 2027 & 60 \\
A5 & Procédure de guichet en mode dégradé, exercée & Récup. & Dir. réseau
d'agences & avril 2027 & 5 \\
A6 & Sensibilisation des agents des 38 agences & Entrée & RSSI & juin 2027
& 10 \\
A7 & Plan de réponse aux incidents testé par un exercice & Récup. & RSSI &
sept. 2027 & 8 \\
A8 & Journalisation centralisée des postes et du partage & Détection &
DSI & juin 2027 & 30 \\
\textbf{Total} & & & & & \textbf{193} \\
\end{tableaucours}
\vspace{-0.8em}
\begin{center}{\small\itshape Tableau~3.6. Plan de traitement de R-AG-01
(décembre 2026)}\end{center}

Deux remarques. L'action A3 n'ajoute aucune mesure~: elle produit la
preuve sans laquelle A2 n'abaisserait rien. Et l'action A8 ne correspond à
aucune étape de la convention~: elle ne réduit pas la cote, mais elle
réduit l'incertitude, puisque la banque saura désormais si elle est
attaquée.

\begin{boxprocedure}
\textbf{Le contenu minimal d'un plan de traitement.} Pour chaque action~:
\begin{itemize}
    \item le risque traité et l'étape du scénario visée~;
    \item le responsable (une personne, pas un service) et l'échéance~;
    \item le coût estimé et la source du financement~;
    \item la \textbf{preuve attendue} qui permettra de la déclarer
          réalisée (procès-verbal de test, extraction, rapport)~;
    \item l'effet attendu sur la cote, selon la convention écrite.
\end{itemize}
Le plan entier porte l'approbation du propriétaire du risque, et sa date
de revue.
\end{boxprocedure}

% ============================================================================
\section{Accepter le Risque Résiduel}
\label{sec:c03-accepter}
% ============================================================================

\subsection{Une décision, pas un constat}

\begin{boxdef}
\termecle{Acceptation du risque}\index{acceptation du risque}~: décision
argumentée, prise par une personne ou une instance qui en a l'autorité, de
ne pas traiter davantage un risque résiduel, pour une durée donnée et à
des conditions données.
\end{boxdef}

L'acceptation est la décision que les registres oublient le plus souvent.
Un risque qui reste dans le registre d'une année sur l'autre, sans
traitement ni signature, est accepté de fait~: personne ne l'a décidé, et
personne n'en répond. C'est la forme la plus discrète du défaut de
gouvernance étudié au Chapitre~\ref{chap:grc}.

\subsection{Le cas de DIGICORE}

Le scénario DC-02, défaillance du prestataire d'hébergement DIGICORE,
était coté $2 \times 3 = 6$ en 2025, avec une vraisemblance sans base
(\NEK{}). Le contrat ne contient pas de clause de droit d'audit~; la banque
ne sait pas comment le prestataire gère ses propres risques.

\begin{boxobjet}{R-DC-02 · Décision d'acceptation (janvier 2027)}
\textbf{Risque~:} défaillance prolongée de DIGICORE, prestataire
d'hébergement du centre de données.

\textbf{Propriétaire~:} le directeur des opérations.

\textbf{Résiduel~:} vraisemblance~2, impact~3, criticité~6. Base de la
vraisemblance~: \NEK{} (aucune revue du prestataire, pas de droit
d'audit). Base de l'impact~: \OUI{} (analyse d'impact sur l'activité).

\textbf{Mesures en place~:} site de secours de Bonabéri~; plan de
continuité approuvé en 2022, jamais testé (aucune réduction retenue).

\textbf{Décision~:} accepté par le comité de sécurité \emph{en l'état des
connaissances}, pour six mois. Conditions~: avenant au contrat introduisant
un droit d'audit et un rapport annuel d'assurance du prestataire (A10)~;
test du plan de continuité avant la revue.

\textbf{Revue~:} juillet 2027. Si l'avenant est refusé, le risque est
présenté au comité des risques du conseil.

\textbf{Signatures~:} directeur des opérations (propriétaire)~; président
du comité de sécurité~; RSSI (avis).
\end{boxobjet}

\begin{boxalert}
\textbf{Violation typique.} Faire signer l'acceptation par le RSSI ou par
le DSI parce qu'ils «~connaissent le sujet~».

\smallskip
\textbf{Règle~:} accepte celui qui a l'autorité et supporte les
conséquences, selon les critères d'acceptation. Le RSSI donne un avis~; il
ne signe pas à la place du propriétaire.
\end{boxalert}

% ============================================================================
\section{Suivre et Revoir}
\label{sec:c03-suivre}
% ============================================================================

\subsection{Le programme 2027, neuf mois plus tard}

Le programme de traitement de 2027 compte les huit actions de R-AG-01 et
deux autres~: A9, conservation et lecture des journaux du pare-feu
d'accès (pour R-DC-03), et A10, avenant au contrat de DIGICORE (pour
R-DC-02). Fin septembre 2027, le comité de sécurité fait le point. Une
action est cotée \OUI{} si elle est réalisée \emph{et} que la preuve
attendue est au dossier, \NON{} si elle n'est pas réalisée, \NEK{} si elle
est déclarée réalisée sans la preuve attendue.

\begin{tableaucours}{cL{5.6cm}cL{4.2cm}}{%
    \tableauHeader{Réf.} &
    \tableauHeader{Action} &
    \tableauHeader{Cote} &
    \tableauHeader{Pièce}
}
A1 & Second facteur & \OUI{} & Extraction de la configuration \\
A2 & Cloisonnement & \OUI{} & Rapport du test A3 \\
A3 & Test d'intrusion interne & \OUI{} & Rapport du prestataire \\
A4 & Sauvegardes hors ligne testées & \OUI{} & Procès-verbaux de
restauration \\
A5 & Mode dégradé exercé & \NEK{} & Procédure rédigée~; aucun exercice
tracé \\
A6 & Sensibilisation & \NEK{} & Listes de présence de 22 agences sur 38 \\
A7 & Exercice de réponse aux incidents & \NON{} & Reporté à 2028 \\
A8 & Journalisation centralisée & \NEK{} & « 20 agences raccordées »,
sans extraction \\
A9 & Journaux du pare-feu d'accès & \OUI{} & Configuration, rapports de
revue mensuels \\
A10 & Avenant DIGICORE & \NON{} & Refus du prestataire \\
\end{tableaucours}
\vspace{-0.8em}
\begin{center}{\small\itshape Tableau~3.7. Avancement du programme de
traitement (septembre 2027)}\end{center}

\begin{boxlecture}{Programme de traitement 2027}
\textbf{Constatation.} 5~\OUI{}, 2~\NON{}, 3~\NEK{} sur 10 actions~:
\[
\mu = 0{,}50~;\quad U = 0{,}30~;\quad
H = -(0{,}5\log_2 0{,}5 + 0{,}2\log_2 0{,}2 + 0{,}3\log_2 0{,}3)
  \approx 0{,}500 + 0{,}464 + 0{,}521 \approx 1{,}49~\text{bit}.
\]

\textbf{Analyse} (confiance élevée). Les actions techniques sont faites et
prouvées~; les actions humaines et organisationnelles (mode dégradé,
sensibilisation, exercice) sont en retard ou sans preuve. Pour R-AG-01,
l'entrée et la propagation sont couvertes~: vraisemblance~2. La
récupération l'est aussi grâce à A4~: impact~3. Le résiduel visé, 6, est
atteint. Mais deux \NEK{} portent sur la récupération (A5) et la détection
(A8)~: le jour d'une attaque, la banque ne sait pas encore si ses agences
sauront servir en mode dégradé. Le refus de DIGICORE (A10) déclenche, selon
la décision de janvier, la présentation de R-DC-02 au conseil.

\medskip
\hfill\lecture{0{,}50}{0{,}30}{1{,}49}
\end{boxlecture}

\subsection{Quand revoir}

Une appréciation se revoit à date fixe (au moins une fois par an), et
aussi chaque fois que le contexte change~: incident significatif, nouveau
système, nouveau prestataire, nouvelle menace signalée par les autorités,
changement de réglementation. Le Chapitre~\ref{chap:grc-pratique} montre
comment des indicateurs de risque déclenchent ces revues sans attendre la
date anniversaire.

% ============================================================================
\section{Le Cadre Camerounais}
\label{sec:c03-cadre}
% ============================================================================

\begin{boxjuris}
\textbf{\loicam{n°~2010/012}, article~26, alinéa~1}~\cite{loi2010012}~:

\medskip
\textit{« Les exploitants des systèmes d'information prennent toutes les
mesures techniques et administratives afin de garantir la sécurité des
services offerts. A cet effet, ils se dotent de systèmes normalisés leur
permettant d'identifier, d'évaluer, de traiter et de gérer continûment les
risques liés à la sécurité des systèmes d'information dans le cadre des
services offerts directement ou indirectement. »}

\medskip
\textbf{Article~29, alinéa~1}~:

\medskip
\textit{« Les exploitants des systèmes d'information ont l'obligation de
conserver les données de connexion et de trafic de leurs systèmes
d'information pendant une période de dix (10) ans. »}

\medskip
\textbf{Portée.} Les quatre verbes de l'article~26 (identifier, évaluer,
traiter, gérer continûment) recouvrent les étapes du processus de
l'ISO~31000~: identification, analyse et évaluation, traitement, suivi et
revue. Une démarche conforme à l'ISO/IEC~27005 est donc une manière de
satisfaire cette obligation, sans être la seule.

\medskip
\textbf{Analyse} (confiance moyenne~: la portée exacte des « données de
connexion et de trafic » est étudiée au Chapitre~\ref{chap:droit}). Des
journaux de pare-feu effacés au bout de trente jours, comme ceux de la
BANQUE SANAGA en 2025, se concilient mal avec l'article~29. L'action A9 ne
traite donc pas seulement un risque~: elle répond aussi à une obligation.
\end{boxjuris}

% ============================================================================
\section{Ce que le Processus demande à chaque Métier}
\label{sec:c03-metiers}
% ============================================================================

\begin{boxmetier}{Gestionnaire des risques cyber}
\textbf{Faire tourner le processus}
\begin{enumerate}
    \item Faire approuver les critères d'acceptation avant toute
          appréciation.
    \item Combiner l'approche par les événements et l'approche par les
          actifs.
    \item Pour chaque risque, obtenir une décision écrite~: un plan, ou une
          acceptation signée.
    \item Suivre le plan par ses preuves, pas par les déclarations.
\end{enumerate}
\end{boxmetier}

\begin{boxmetier}{RSSI}
\textbf{Un plan qui produit des preuves}

Chaque action du plan prévoit la preuve qui la clôturera. Le RSSI inscrit
au plan les tests eux-mêmes (comme A3), car une mesure non testée
n'abaisse rien, et refuse de déclarer réalisée une action sans sa pièce.
\end{boxmetier}

\begin{boxmetier}{Consultant en mise en œuvre de SMSI}
\textbf{Livrer un processus, pas un registre}

Le consultant qui accompagne une organisation laisse derrière lui des
critères approuvés, une méthode écrite, un plan avec des responsables et
une première revue datée. Un registre bien présenté, livré seul, sera
obsolète à la première réorganisation.
\end{boxmetier}

\begin{boxmetier}{Auditeur en cybersécurité}
\textbf{Auditer les décisions}

L'auditeur prélève des risques du registre et demande, pour chacun, la
décision~: le plan approuvé ou l'acceptation signée. Il vérifie que le
signataire avait l'autorité selon les critères, et qu'aucune acceptation ne
repose sur une cote sans base, sauf acceptation datée en l'état des
connaissances.
\end{boxmetier}

\separateur

\begin{boxresume}
\textbf{Ce qu'il faut retenir de ce chapitre~:}
\begin{itemize}
    \item \textbf{L'ISO~31000} donne des principes, un cadre et un
          processus~; la communication et le suivi encadrent toutes les
          étapes.
    \item \textbf{Le contexte} fixe le périmètre, les parties intéressées
          et deux familles de critères~: pour apprécier, et pour accepter.
    \item \textbf{Deux approches d'identification}~: par les événements
          pour savoir ce qui compte, par les actifs pour vérifier que rien
          n'échappe.
    \item \textbf{Quatre familles de traitement}~: réduire, partager,
          éviter, maintenir~; l'assurance vient après la réduction.
    \item \textbf{Un plan de traitement} donne, pour chaque action, un
          responsable, une échéance, un coût et la preuve attendue.
    \item \textbf{L'acceptation} est une décision écrite, par l'instance
          compétente, à durée limitée~; sur une cote \NEK{}, elle se fait
          « en l'état des connaissances », avec une date de revue.
    \item \textbf{Le suivi} cote les actions par leurs preuves~: en
          septembre 2027, $\mu = 0{,}50$, $U = 0{,}30$ à la BANQUE SANAGA.
\end{itemize}
\end{boxresume}

\begin{boxinfo}
\textbf{Points de contrôle.} Vous avez compris ce chapitre si vous savez~:
\begin{itemize}
    \item placer une activité dans le processus de l'ISO~31000~;
    \item dire ce que l'approche par les actifs apporte à l'approche par les
          événements, sur un exemple~;
    \item rédiger une ligne de plan de traitement complète~;
    \item dire qui peut accepter un résiduel donné, et à quelles
          conditions.
\end{itemize}
\end{boxinfo}

% ============================================================================
\section*{Exercices}
% ============================================================================

\begin{boxexo}[title={\ding{45}\quad Exercice~3.1. Où en est-on dans le processus~? \niveaubase}]
Placez chacune des activités suivantes dans le processus de l'ISO~31000
(Figure~\ref{fig:c03-processus}).
\begin{enumerate}[(a)]
    \item Le comité des risques fixe qui peut accepter un résiduel de 5 à 8.
    \item L'inventaire des postes des 38 agences.
    \item La RSSI cote la vraisemblance de la fraude au virement.
    \item Le comité compare les criticités aux critères et classe les
          risques.
    \item La DSI déploie le second facteur.
    \item Le comité de sécurité fait le point sur le programme 2027.
    \item La direction du réseau d'agences participe à l'atelier des
          scénarios.
    \item Le rapport annuel des risques est transmis au conseil.
\end{enumerate}
\end{boxexo}

\begin{boxexo}[title={\ding{45}\quad Exercice~3.2. Choisir un traitement \niveaumoyen}]
Pour chacun des risques suivants de LiZbiZ (appétence~: 9 et plus
intolérable, 5 à 8 traité dans l'année, 1 à 4 accepté), proposez une
famille de traitement et une mesure, et justifiez.
\begin{enumerate}[(a)]
    \item R1, ajustement frauduleux du stock, résiduel 12.
    \item R4, arrêt prolongé du centre de Garoua, résiduel 4, base de
          probabilité \OUI{}.
    \item R6, fuite de données des clients, résiduel 3, base de probabilité
          \NEK{} (Exercice~2.4).
    \item R8, rançongiciel, résiduel 12 après nouvelle cote.
\end{enumerate}
Indiquez, pour (d), pourquoi une assurance seule ne suffit pas.
\end{boxexo}

\begin{boxexo}[title={\ding{45}\quad Exercice~3.3. Une acceptation contestable \niveaumoyen}]
En février 2027, le dossier de R-DC-06 (déni de service des services en
ligne, coté $1 \times 3 = 3$ en 2025) contient la décision suivante~:
« Risque accepté. Signé~: le DSI. » Aucune date de revue n'est indiquée.
\begin{enumerate}[(a)]
    \item Relevez tous les défauts de cette décision au regard des critères
          d'acceptation (Tableau~3.3).
    \item La base de la vraisemblance était cotée \NON{} (quatre attaques
          signalées en 2025). Proposez une nouvelle vraisemblance, recalculez
          la criticité et dites qui peut désormais accepter le risque.
    \item Rédigez la décision corrigée sur le modèle de R-DC-02.
\end{enumerate}
\end{boxexo}

\begin{boxexo}[title={\ding{45}\quad Exercice~3.4. Le plan de traitement de R8 \niveauexpert}]
LiZbiZ décide de traiter R8 (rançongiciel, probabilité~3, impact~4,
criticité~12, aucune mesure testée). L'audit des quatre centres a relevé
des comptes par défaut, des services en clair et des flux non autorisés.
\begin{enumerate}[(a)]
    \item Décomposez le scénario en étapes (entrée, propagation,
          récupération) et reliez chaque constat de l'audit à une étape.
    \item Rédigez un plan de traitement de six actions au plus, au format du
          Tableau~3.6, en indiquant pour chacune la preuve attendue.
    \item Calculez le résiduel visé avec la convention de réduction, en
          partant d'un inhérent de $4 \times 4$, et dites qui pourra
          l'accepter selon des critères identiques à ceux de la BANQUE
          SANAGA.
    \item Un an plus tard, quatre actions sont prouvées, une est déclarée
          faite sans preuve, une n'est pas faite. Calculez $\mu$, $U$ et
          $H$, et dites si le résiduel visé est atteint.
\end{enumerate}
\end{boxexo}

\bigskip

\begin{boxinfo}
\textbf{Pour aller plus loin.}
\begin{itemize}
    \item L'ISO~31000:2018~\cite{iso31000_2018}, pour les principes, le
          cadre et le processus.
    \item L'ISO/IEC~27005:2022~\cite{iso27005_2022}, pour l'application à
          la sécurité de l'information.
    \item L'ISO/IEC~27001:2022~\cite{iso27001_2022}, article~6.1.3, sur le
          traitement des risques et la déclaration d'applicabilité.
\end{itemize}
\end{boxinfo}

\begin{boxinfo}
\textbf{Transition.} Le processus dit quoi faire~; il ne dit pas comment
imaginer les scénarios. L'atelier de décembre 2026 a trouvé quatre
scénarios stratégiques~: comment les construire avec méthode, à partir des
sources de risque et de leurs objectifs, puis les détailler en modes
opératoires~? C'est l'objet de la méthode EBIOS Risk Manager, au
Chapitre~\ref{chap:ebios}.
\end{boxinfo}
